\section{Introduction}
\label{sec:intro}

XLA, the compiler behind Google's TPUs, must decide how to arrange each neural network's tensors in
memory. There are many valid arrangements, called layout configurations, and they run at measurably
different speeds on real hardware. Timing every candidate is slow and the search space is too large to
try exhaustively, so a learned cost model ranks the candidates from the program's graph instead of
keeping the compiler in the loop~\cite{phothilimthana2023tpugraphs,kaufman2021learned}. Layout matters
for a physical reason: the TPU's matrix multiply unit (MXU) works on fixed-size tiles, so the way data
is laid out changes how efficiently the MXU is used~\cite{jouppi2020domainspecific}.

Progress on these models is reported almost entirely as architecture comparisons: model X ranks
configurations better than model Y. GNNs over computation graphs have become the dominant
approach~\cite{phothilimthana2023tpugraphs,kaufman2021learned,deng2024polynormer}. Such comparisons
cannot show whether a gain is real unless they control for seed noise, model capacity and baseline
strength. In an audit of 20 papers that make these comparisons, none runs a significance test on its
headline claim, and 17 of 20 fail at least three of five rigor axes (Section~\ref{sec:audit}).

We therefore ask a different question: how much of a configuration's speed is knowable from the static
graph the predictor sees, and how much lives only in the compiler's dynamic tiling search? We answer it
at three levels:
\begin{enumerate}
\item Does the static graph carry the performance information, and does that depend on the workload?
\item Where, physically, does it fall short?
\item Can current evaluation practice detect effects of this size?
\end{enumerate}

At each level the answer turns on baseline strength. First, our tested alterations (hierarchy and
contention edges) to graph structure provided no detectable gain over a flat matched model in any of
the 4 collections once the baseline is tuned: on xla:default, an apparent $+0.165$ $\tau$ at a
shared learning rate falls to $+0.002$ ($p = .88$). Second, convolution graphs are $2.8\times$ more
layout-sensitive than other graphs in measured runtimes, but a per-node MXU feature that narrows
convolution slowdown from 1.42 to 1.12 at the shared rate leaves a gain below detection ($-0.03$, $p
= .17$) against a tuned baseline. Third, a result that can pass seed, parameter-matching and
significance checks can still be a baseline artefact, as our own $+0.165$ was. All results are on
TPUGraphs' validation graphs with runtimes Google measured; we ran no TPU.

Our contributions are:
\begin{itemize}
\item A controlled and power-analysed protocol for comparing learned cost models, applied to all four
  TPUGraphs layout collections.
\item A model-free measurement that localises the information gap to convolution on the MXU.
\item A pre-registered test of whether per-node MXU geometry closes that gap, with a scramble control, a
  tuned-baseline check and an oracle ceiling.
\item An audit of evaluation practice across papers in the field along with concrete recommendations
  (Section~\ref{sec:recs}), the released protocol and the power-analysis tool.
\end{itemize}
